\subsection{Executing an original call}
A dispatch freezes the complete model input, source audio range, target
frontier and execution contract. Each native FP32 worker executes the
original inference, selection and feedback-update code. Its early ASR
observation can supply a typed provisional view while original correction
continues; its complete result follows the original C/U path.

All native policies use the same work-conserving dispatcher. It first
considers eligible calls from actual selected inputs across sessions,
then conditional calls, and rotates among eligible sessions. It registers
and sends one dispatch before reconsidering the next free slot.
Conservative can therefore use both GPUs for distinct sessions under the
same resource budget. The policy differences concern which outstanding
calls survive revisions, as evaluated in
Section~\ref{sec:eval}.

\subsection{Logical invalidation and physical ownership}
Dispatch and receipt identities bind the session, Owner namespace,
monotonic job number and physical worker incarnation, following established
RPC identity principles~\cite{birrell1984rpc}. Exact duplicates and fenced
messages are rejected locally. They neither modify selected state nor
release a newer job's slot. Contradictory contents from a currently owned
call freeze further writes.

A retention decision does not terminate GPU execution. In the native
runtime, a known-input mismatch or covered target can request a cooperative
exit at the checkpoint before learned correction. A worker either returns
a complete original result or acknowledges cancellation without a partial
feedback state. Cleanup and device synchronization precede its terminal
receipt. If completion wins the race, the returned result still passes
the original input and selection checks. UIR suppresses these early-exit
requests while keeping the same end-use checks.

Per-incarnation channels bind the message stream to its physical producer.
After an unexpected channel close, complete received frames are validated
before the channel is sealed. Replacement requires confirmed process
exit, released CUDA ownership and joined receivers. In the batched core
instance, a completed row can be used before other rows finish, but the
batch lease ends only after every terminal row and device synchronization.
Neither a logical fence nor a masked row is charged as freed computation.

\subsection{Durable selection and recovery}
An exclusive writer records the full parent and next states, original
correction and returned IDs, owned job and relevant epochs. Framing,
checksums and chaining bind the record; synchronization precedes
installation and publication, following conventional write-ahead
rules~\cite{mohan1992aries}. A successful callback returns the receipt for
that record. An uncertain append or installation outcome freezes writes
until recovery establishes the verified prefix.

\label{subsec:native_resilience}
For an Owner-group failure, the Supervisor confirms the end of old
physical ownership, advances the Owner namespace, and replays the verified
durable prefix to the surviving Consumer before starting the replacement
Owner. Only an incomplete final frame may be discarded; complete
corruption fails verification. Replay preserves the original fact and
job identities. Workers cold-start and continue from the recovered
selected input, using the original ingress clocks.

The sync--install distinction has an observed counterpart. In two repeats
of an earlier Fullcall fault case, journal sequence 5 was synchronized when
the original Owner was terminated before installation. During recovery
to the next Owner epoch, the surviving Consumer replayed that same fact
with its original job identity. For these two facts, the counts of
durable records, original-Owner installations, and Consumer uses were
$2/0/2$. This instantiates the sync--install gap permitted by
Proposition~\ref{prop:refinement}: replay used the existing edge without
adding a neural transition.

Twenty-four further fault cells on one development source exercise six
conditions under both policies. No cell installs illegally or accepts a
rejected mutation. Recovery cost splits by what must be rebuilt: worker
faults resume in 1.8--2.3~s against a surviving model, while Owner-group
faults pay 64.6--108.1~s to rebuild the worker group. No fault cell meets
its deadlines, so these trials establish the absence of an illegal
transition and a bound on recovery rather than continued timeliness. One
condition prices retention directly: when the Owner dies after the
durable install, retaining needs 101.8--104.4~s and 26.5--27.2~kJ against
83.9--87.1~s and 22.4--24.2~kJ, because more outstanding work survives to
be re-driven. Retention moves cost onto the recovery path instead of
paying it in steady state.

An independent audit of that earlier runtime also bound 3,976 selected
and 9,068 provisional applications across 36 traces, and verified
preservation at all 90 compatible-ahead applications.

\phantomsection\label{prop:continued_service}
After finitely many supported failures, eventual progress requires readable
storage, terminating original calls, eventual worker release and loading,
and fair Owner/Consumer scheduling. Loading and replay remain part of
post-fault service time; destroying a model process incurs its actual
reconstruction cost even when an acquired result survives.

\subsection{Keeping control history bounded}
Active calls, ready results, delivered views and unacknowledged
publications pin their required evidence. Acknowledgment releases
publication references, not arbitrary computation dependencies. Monotonic
issue high watermarks prevent old identities from being reissued after
collection. Unreferenced terminal and diagnostic windows are bounded and
streamed to disk. These are bounds on control history; audio, cumulative
text and model state retain their separately measured costs.
